Fix \(n\geq256\), \(0<\Delta\leq\Delta_0\), and \(C\in\mathcal C_{n,\Delta}(\alpha)\).  By \(\mathrm{thm:single\text{-}joint\text{-}pair}\), the two laws in \(\mathfrak H_{n,\Delta}\) are members of \(\mathcal P_L\).  Define
\[
 Q_j=\mathcal L_{P_{j,n,\Delta}}(Z_\Delta^n),
 \qquad
 \vartheta_j=\tau_{P_{j,n,\Delta}}(b),
 \qquad j\in\{0,1\}.
\]
The same theorem, with \(\kappa=c_L/4\), gives
\[
 \operatorname{TV}(Q_0,Q_1)\leq\frac14
 \quad\text{and}\quad
 \delta:=|\vartheta_1-\vartheta_0|
 =\kappa(h_n+\Delta).
\]
Because \(\rho_{n,\Delta}=\max\{h_n,\Delta\}\), each summand is nonnegative and therefore
\[
 h_n+\Delta\geq\rho_{n,\Delta},
 \qquad
 \delta\geq\kappa\rho_{n,\Delta}.
\]

By \(C\in\mathcal C_{n,\Delta}(\alpha)\) and \(P_{j,n,\Delta}\in\mathcal P_L\), the honesty condition in \(\mathrm{def:honest\text{-}class}\) yields
\[
 Q_0\{\vartheta_0\in C\}\geq1-\alpha,
 \qquad
 Q_1\{\vartheta_1\in C\}\geq1-\alpha.
\]
The events are measurable because \(C\) is a \(\sigma(Z_\Delta^n)\)-measurable interval-valued map.  For \(A_1=\{\vartheta_1\in C\}\), the defining event inequality for total variation gives
\[
 Q_0(A_1)\geq Q_1(A_1)-\operatorname{TV}(Q_0,Q_1).
\]
Consequently, the elementary inequality \(Q(A\cap B)\geq Q(A)+Q(B)-1\) implies
\[
 \begin{aligned}
 Q_0\{\vartheta_0\in C,\ \vartheta_1\in C\}
 &\geq Q_0\{\vartheta_0\in C\}+Q_0(A_1)-1\\
 &\geq1-2\alpha-\operatorname{TV}(Q_0,Q_1)\\
 &\geq\frac34-2\alpha.
 \end{aligned}
\]
The last quantity is positive because \(\alpha\in(0,1/4)\).  On the displayed intersection, the real interval \(C\) contains both \(\vartheta_0\) and \(\vartheta_1\), so \(\operatorname{length}(C)\geq\delta\).  Nonnegativity of interval length and the membership \(P_{0,n,\Delta}\in\mathcal P_L\) now give
\[
 \begin{aligned}
 \sup_{P\in\mathcal P_L}\mathbb E_P[\operatorname{length}(C)]
 &\geq \mathbb E_{P_{0,n,\Delta}}[\operatorname{length}(C)]\\
 &\geq \delta\,Q_0\{\vartheta_0\in C,\ \vartheta_1\in C\}\\
 &\geq \frac{c_L}{4}\left(\frac34-2\alpha\right)\rho_{n,\Delta}.
 \end{aligned}
\]
Thus the asserted lower bound holds with
\[
 c_{L,\alpha,\Delta_0}
 :=\frac{c_L}{4}\left(\frac34-2\alpha\right)>0.
\]
This constant depends on \(L\) through \(c_L\) and on \(\alpha\), and is independent of \(n\), \(\Delta\), \(c_0\), and \(\Delta_0\); this is stronger than the permitted dependence.  Taking the infimum over \(C\in\mathcal C_{n,\Delta}(\alpha)\) and using \(\mathrm{def:honest\text{-}width}\) proves the minimax lower bound for \(\mathfrak W_{n,\Delta}(\alpha)\).

For the matching upper bound, \(\mathrm{thm:total\text{-}honest\text{-}interval}\), together with \(\mathrm{def:total\text{-}honest\text{-}interval}\) and \(\mathrm{def:honest\text{-}class}\), shows that the explicit closed bounded released-sample-measurable interval \(C_{n,\Delta,\alpha}\) belongs to \(\mathcal C_{n,\Delta}(\alpha)\) and satisfies
\[
 \sup_{P\in\mathcal P_L}
 \mathbb E_P[\operatorname{length}(C_{n,\Delta,\alpha})]
 \leq C_{L,c_0,\alpha,\Delta_0}\rho_{n,\Delta}.
\]
The infimum defining \(\mathfrak W_{n,\Delta}(\alpha)\) is no larger than the value at this interval.  Combining the two inequalities proves \(\mathfrak W_{n,\Delta}(\alpha)\asymp\rho_{n,\Delta}\), uniformly over all stated \(n\) and \(\Delta\).